\chapter{Project Requirements}

This chapter describes the comprehensive functional and non-functional requirements for QCanvas, followed by use-cases, external interface specifications, and system features that define the complete system behavior.

\section{Functional Requirements}

\subsection{Conversion Module}
\begin{itemize}
    \item FR-C1: The system shall parse Cirq, Qiskit, and PennyLane circuits using AST-based parsing with runtime fallback.
    \item FR-C2: The system shall generate OpenQASM 3.0 code conforming to the supported subset per project scope.
    \item FR-C3: The system shall handle measurement and reset instructions correctly.
    \item FR-C4: The system shall expose conversion statistics including qubits, depth, and gate counts.
    \item FR-C5: The system shall validate generated OpenQASM 3.0 against the supported subset and provide detailed diagnostics.
    \item FR-C6: The system shall support gate mapping with parameter and broadcast handling.
\end{itemize}

\subsection{Simulation Module}
\begin{itemize}
    \item FR-S1: The system shall support statevector, density matrix, and stabilizer simulation backends.
    \item FR-S2: The system shall allow configurable shots (1–10,000) and return counts and metadata.
    \item FR-S3: The system shall cache identical executions to improve performance.
    \item FR-S4: The system shall allow seed configuration for deterministic runs.
    \item FR-S5: The system shall provide backend capability discovery (max qubits, supported gates).
    \item FR-S6: The system shall return results with counts, execution time, and optional state snapshots.
\end{itemize}

\subsection{Hybrid CPU–QPU Module}
\begin{itemize}
    \item FR-H1: The system shall separate classical orchestration (CPU) from quantum circuit execution (QPU/Simulator).
    \item FR-H2: The system shall implement data flow: Source $\rightarrow$ QASM $\rightarrow$ Execution $\rightarrow$ Results.
    \item FR-H3: The system shall provide WebSocket progress updates for long-running operations.
    \item FR-H4: The system shall enforce simulation-first execution with optional pluggable endpoints.
    \item FR-H5: The system shall provide resource-aware execution with limits and timeouts.
\end{itemize}

\subsection{Web IDE Module}
\begin{itemize}
    \item FR-W1: The system shall provide code editing with syntax highlighting and diagnostics.
    \item FR-W2: The system shall display circuit visualization and results panels.
    \item FR-W3: The system shall support multiple example projects and templates.
    \item FR-W4: The system shall provide keyboard shortcuts for frequent actions.
    \item FR-W5: The system shall support multi-file projects with workspace management.
\end{itemize}

\subsection{API and Realtime}
\begin{itemize}
    \item FR-A1: The system shall expose REST endpoints for conversion and validation.
    \item FR-A2: The system shall expose WebSocket endpoints for progress streaming.
    \item FR-A3: The system shall provide structured error payloads with codes and hints.
    \item FR-A4: The system shall provide discovery endpoints for frameworks and backends.
    \item FR-A5: The system shall support batch conversion and simulation operations.
\end{itemize}

\subsection{Validation}
\begin{itemize}
    \item FR-V1: The system shall validate OpenQASM 3.0 output against the allowed subset.
    \item FR-V2: The system shall suggest nearest supported constructs for incompatibilities.
    \item FR-V3: The system shall emit precise diagnostics with line/column, code, and remediation.
    \item FR-V4: The system shall check for excluded features (timing, pulse-level) with clear guidance.
\end{itemize}

\section{Non-Functional Requirements}

\subsection{Reliability}
\begin{itemize}
    \item NFR-R1: The system shall provide graceful degradation without crashing the IDE.
    \item NFR-R2: Core services shall recover from transient faults without full termination.
    \item NFR-R3: The system shall maintain persistent logs for failures and audit trails.
    \item NFR-R4: The system shall target 99.9\% uptime for core services.
\end{itemize}

\subsection{Usability}
\begin{itemize}
    \item NFR-U1: The system shall provide responsive UI with clear error messages.
    \item NFR-U2: The system shall support keyboard-first workflow with minimal steps to run.
    \item NFR-U3: The system shall provide accessible design with high-contrast theme and screen-reader support.
    \item NFR-U4: The system shall allow compilation or execution within 3 user interactions from editor.
\end{itemize}

\subsection{Performance}
\begin{itemize}
    \item NFR-P1: API response time shall be less than 2 seconds for typical conversions.
    \item NFR-P2: Simulation time shall be less than 30 seconds for circuits up to 20 qubits (statevector).
    \item NFR-P3: IDE interactions shall stay under 100ms perceived latency for common actions.
    \item NFR-P4: WebSocket updates shall have less than 100ms latency.
\end{itemize}

\subsection{Security}
\begin{itemize}
    \item NFR-S1: The system shall validate all inputs and enforce rate limiting on endpoints.
    \item NFR-S2: The system shall use secure WebSocket connections in production.
    \item NFR-S3: The system shall implement CORS configuration and CSRF protections.
    \item NFR-S4: The system shall sanitize user-supplied code and logs to prevent injection.
\end{itemize}

\clearpage
\section{External Interface Requirements}

\subsection{User Interfaces}

\subsubsection{Web IDE Interface}
\begin{itemize}
    \item Monaco-based code editor with syntax highlighting for Qiskit, Cirq, PennyLane, and OpenQASM 3.0
    \item Multi-panel layout: Editor, Results, Console, Logs, Statistics
    \item File tree navigation with create, rename, delete operations
    \item Theme switcher (dark/light modes)
    \item Responsive design for desktop and tablet devices
\end{itemize}

\subsubsection{Results Visualization}
\begin{itemize}
    \item Tabbed results display: Console, Output, Histogram, OpenQASM, Errors, Stats
    \item Interactive measurement histograms using Chart.js
    \item Real-time execution status updates
    \item Export functionality for results (JSON, CSV)
\end{itemize}

\subsection{Hardware Interfaces}

The system does not interface with physical quantum hardware. All execution is simulation-based through software backends.

\subsection{Software Interfaces}

\subsubsection{Quantum Framework APIs}
\begin{itemize}
    \item Qiskit API (version 0.44+): Circuit construction and manipulation
    \item Cirq API (version 1.0+): Circuit operations and gate definitions
    \item PennyLane API (version 0.33+): Quantum node processing
\end{itemize}

\subsubsection{Database Interface}
\begin{itemize}
    \item PostgreSQL (future): User data, projects, circuit storage
    \item Redis (future): Session caching and result caching
\end{itemize}

\subsubsection{External Services}
\begin{itemize}
    \item No external quantum hardware services in current scope
    \item Future: Pluggable QPU endpoints with same API contract
\end{itemize}

\subsection{Communication Interfaces}

\subsubsection{REST API}
\begin{itemize}
    \item Protocol: HTTPS in production, HTTP in development
    \item Format: JSON request/response bodies
    \item Authentication: JWT tokens (future implementation)
    \item Base URL: /api/v1/
\end{itemize}

\subsubsection{WebSocket Interface}
\begin{itemize}
    \item Protocol: WSS in production, WS in development
    \item Format: JSON message payloads
    \item Endpoint: /ws
    \item Use cases: Progress updates, real-time notifications
\end{itemize}

\clearpage
\section{Use-Cases}

The following use cases describe the core interactions between users and the QCanvas system.

\subsection{UC-01: Launch Web IDE}
\textbf{Primary Actor:} User \\
\textbf{Precondition:} Browser access to QCanvas platform \\
\textbf{Main Flow:}
\begin{enumerate}
    \item User navigates to QCanvas web application
    \item System loads Monaco editor and initializes workspace
    \item System displays example library and framework selector
\end{enumerate}
\textbf{Postcondition:} IDE ready for circuit input

\subsection{UC-02: Write Circuit Code}
\textbf{Primary Actor:} User \\
\textbf{Precondition:} Web IDE loaded successfully \\
\textbf{Main Flow:}
\begin{enumerate}
    \item User selects framework (Cirq, Qiskit, or PennyLane)
    \item User writes quantum circuit code in editor
    \item System provides syntax highlighting and basic validation
\end{enumerate}
\textbf{Postcondition:} Circuit code ready for conversion

\subsection{UC-03: Convert to OpenQASM}
\textbf{Primary Actor:} User \\
\textbf{Precondition:} Valid circuit code exists in editor \\
\textbf{Main Flow:}
\begin{enumerate}
    \item User clicks "Convert to QASM" button
    \item QCanvas (CPU) parses source framework code via AST or runtime
    \item QCanvas generates OpenQASM 3.0 intermediate representation
    \item System validates generated QASM against supported subset
    \item System displays converted QASM code and statistics
\end{enumerate}
\textbf{Postcondition:} Valid OpenQASM 3.0 code available

\subsection{UC-04: Execute Circuit on Simulator}
\textbf{Primary Actor:} User \\
\textbf{Precondition:} Valid OpenQASM 3.0 code exists \\
\textbf{Main Flow:}
\begin{enumerate}
    \item User selects simulation backend (statevector, density matrix, or stabilizer)
    \item User configures execution parameters (shots, seed)
    \item QCanvas (CPU) submits job to QSim (QPU/Simulator)
    \item QSim executes circuit gates and performs measurements
    \item System streams progress updates via WebSocket
    \item QSim returns measurement results to QCanvas
\end{enumerate}
\textbf{Postcondition:} Simulation results generated and available

\subsection{UC-05: Visualize Results}
\textbf{Primary Actor:} User \\
\textbf{Precondition:} Simulation execution completed successfully \\
\textbf{Main Flow:}
\begin{enumerate}
    \item System renders measurement histogram using Chart.js
    \item System displays execution statistics (time, backend info)
    \item System shows optional state snapshots (statevector/density matrix)
\end{enumerate}
\textbf{Postcondition:} Results visualized in IDE

\subsection{UC-06: Handle Validation Errors}
\textbf{Primary Actor:} User \\
\textbf{Precondition:} Validation or execution error occurred \\
\textbf{Main Flow:}
\begin{enumerate}
    \item System displays error message with line/column information
    \item System highlights problematic code in editor
    \item System suggests fixes or identifies unsupported features
    \item User corrects code based on feedback
\end{enumerate}
\textbf{Postcondition:} User informed of issues and guided to resolution

\subsection{UC-07: Export OpenQASM}
\textbf{Primary Actor:} User \\
\textbf{Precondition:} OpenQASM 3.0 code successfully generated \\
\textbf{Main Flow:}
\begin{enumerate}
    \item User clicks export button
    \item System packages QASM code with metadata
    \item System downloads .qasm file to user's machine
\end{enumerate}
\textbf{Postcondition:} QASM file saved locally

\subsection{UC-08: Load Example Circuit}
\textbf{Primary Actor:} User \\
\textbf{Precondition:} IDE loaded and example library available \\
\textbf{Main Flow:}
\begin{enumerate}
    \item User browses curated example library
    \item User selects example (e.g., Bell State, GHZ, Grover's Algorithm)
    \item System loads example code into editor with annotations
\end{enumerate}
\textbf{Postcondition:} Example ready to convert and execute

\subsection{UC-09: Batch Circuit Conversion}
\textbf{Primary Actor:} Researcher \\
\textbf{Precondition:} Multiple circuit files prepared \\
\textbf{Main Flow:}
\begin{enumerate}
    \item Researcher submits batch conversion request via API
    \item QCanvas processes each circuit sequentially or in parallel
    \item System aggregates conversion results and statistics
    \item System returns batch results with individual diagnostics
\end{enumerate}
\textbf{Postcondition:} All circuits converted with detailed reports

\subsection{UC-10: Compare Framework Outputs}
\textbf{Primary Actor:} Researcher \\
\textbf{Precondition:} Same circuit written in multiple frameworks \\
\textbf{Main Flow:}
\begin{enumerate}
    \item Researcher converts all framework versions to OpenQASM
    \item System executes all circuits on same backend with same parameters
    \item System compares measurement distributions and statistics
    \item System generates comparison report
\end{enumerate}
\textbf{Postcondition:} Cross-framework comparison data available

\subsection{UC-11: Query Backend Capabilities}
\textbf{Primary Actor:} Developer \\
\textbf{Precondition:} System running and API accessible \\
\textbf{Main Flow:}
\begin{enumerate}
    \item Developer queries backend info endpoint
    \item System returns maximum qubits supported
    \item System lists supported gates and operations
    \item System provides backend-specific constraints
\end{enumerate}
\textbf{Postcondition:} Backend capabilities known for optimization

\subsection{UC-12: Hybrid CPU–QPU Execution}
\textbf{Primary Actor:} User/System \\
\textbf{Precondition:} Valid circuit ready for execution \\
\textbf{Main Flow:}
\begin{enumerate}
    \item QCanvas (CPU) handles circuit compilation and validation
    \item QCanvas prepares execution parameters and context
    \item QCanvas submits compiled circuit to QSim (QPU/Simulator)
    \item QSim (QPU) performs pure quantum gate execution
    \item QSim returns raw measurement results
    \item QCanvas (CPU) aggregates and post-processes results
    \item System delivers final results to user via API
\end{enumerate}
\textbf{Postcondition:} Hybrid execution complete with clear CPU–QPU separation

\clearpage
\section{Use Cases Diagram}
This diagram outlines the primary actors and their interactions with QCanvas across core flows like conversion, simulation, and results visualization.
\begin{figure}[H]
    \centering
    \includegraphics[width=\linewidth]{Use Case Diagram.png}
    \caption{Use Case Diagram}
\end{figure}

\clearpage
\section{Domain Model Diagram}
This diagram shows the core entities (e.g., Circuit, ConversionRequest, SimulationJob, SimulationResult, User, Project) and key relationships that support conversion and simulation.
\begin{figure}[H]
    \centering
    \includegraphics[width=\linewidth]{Domain Model.png}
    \caption{Domain Model}
\end{figure}
